
\documentclass[12pt,a4paper]{article}

\usepackage[utf8]{inputenc}
\usepackage[T5]{fontenc}
\usepackage[vietnamese]{babel}
\usepackage{amsmath}
\usepackage{amssymb}
\usepackage{geometry}
\usepackage{tikz}
\usepackage{longtable}
\usepackage{hyperref}

\geometry{margin=2.5cm}

\title{Lộ Trình Xây Dựng Chatbot Da Liễu \\ Từ HAM10000 (7 Bệnh) Đến Hàng Trăm Bệnh}
\author{System Design Document}
\date{\today}

\begin{document}

\maketitle

\tableofcontents
\newpage

\section{Mục tiêu hệ thống}

Mục tiêu là xây dựng một chatbot da liễu có khả năng:

\begin{itemize}
\item Phân tích ảnh tổn thương da.
\item Trả lời câu hỏi bằng ngôn ngữ tự nhiên.
\item Giải thích lý do đưa ra dự đoán.
\item Truy xuất tri thức y khoa.
\item Mở rộng từ 7 bệnh HAM10000 lên hàng trăm bệnh da liễu.
\end{itemize}

\section{Tại sao không chỉ dùng CNN?}

CNN có thể dự đoán bệnh từ ảnh:

\begin{verbatim}
Image
  |
  v
CNN
  |
  v
MEL = 0.85
BCC = 0.10
NV  = 0.05
\end{verbatim}

Tuy nhiên CNN không giải thích được:

\begin{itemize}
\item Vì sao mô hình chọn MEL?
\item Triệu chứng nào được sử dụng?
\item Tài liệu y khoa nào hỗ trợ kết quả?
\end{itemize}

Do đó cần kết hợp nhiều thành phần khác ngoài CNN.

\section{Kiến thức nền tảng}

\subsection{CNN}

CNN (Convolutional Neural Network) là mô hình xử lý ảnh.

Pipeline:

\begin{verbatim}
Image
  |
  v
Feature Maps
  |
  v
Classifier
  |
  v
Disease Probabilities
\end{verbatim}

Kiến thức cần biết:

\begin{itemize}
\item Convolution
\item Pooling
\item Transfer Learning
\item Fine-tuning
\item ResNet
\item EfficientNet
\end{itemize}

\subsection{Embedding}

Embedding biến văn bản thành vector số.

Ví dụ:

\begin{verbatim}
"many colors"
        |
        v
 [0.21, -0.34, ...]
\end{verbatim}

Mục đích:

\begin{itemize}
\item So sánh ngữ nghĩa.
\item Truy xuất tri thức.
\item Semantic Search.
\end{itemize}

\subsection{Sentence Transformer}

Ví dụ:

\begin{verbatim}
all-MiniLM-L6-v2
\end{verbatim}

Vai trò:

\begin{verbatim}
Text
 |
 v
Embedding Vector
\end{verbatim}

\section{Vector Database và FAISS}

\subsection{Vấn đề}

Giả sử có:

\begin{itemize}
\item 100000 đoạn tài liệu y khoa.
\item Hàng trăm bệnh.
\end{itemize}

Không thể so sánh thủ công với mọi đoạn văn.

\subsection{FAISS}

FAISS là thư viện tìm kiếm vector gần nhất.

Pipeline:

\begin{verbatim}
Query
  |
  v
Embedding
  |
  v
FAISS Search
  |
  v
Top-K Chunks
\end{verbatim}

\section{Knowledge Graph}

\subsection{Knowledge Graph là gì?}

Knowledge Graph lưu tri thức dưới dạng đồ thị.

Ví dụ:

\begin{verbatim}
MEL
 |
 +-- asymmetry
 |
 +-- border_irregularity
 |
 +-- color_variation
\end{verbatim}

\subsection{Các loại node}

\begin{itemize}
\item Disease
\item Feature
\item Symptom
\item Risk Factor
\item Diagnosis
\item Treatment
\end{itemize}

\subsection{Các loại quan hệ}

\begin{itemize}
\item HAS\_FEATURE
\item HAS\_SYMPTOM
\item HAS\_RISK
\item DIAGNOSED\_BY
\item TREATED\_BY
\end{itemize}

\section{Semantic Feature Retrieval}

\subsection{Semantic Retrieval là gì?}

Mục tiêu:

\begin{verbatim}
uneven edges
      |
      v
border_irregularity
\end{verbatim}

Hoặc:

\begin{verbatim}
many colors
      |
      v
color_variation
\end{verbatim}

\subsection{Tại sao cần?}

Nếu dùng rule-based:

\begin{verbatim}
FEATURE_MAP = {...}
\end{verbatim}

thì khi mở rộng lên hàng trăm bệnh sẽ rất khó bảo trì.

Semantic Retrieval cho phép tìm feature bằng ý nghĩa thay vì từ khóa cố định.

\subsection{Pipeline}

\begin{verbatim}
Query
  |
  v
Embedding
  |
  v
Feature Index
  |
  v
Relevant Features
\end{verbatim}

\section{Disease Knowledge Retrieval}

Mục tiêu là truy xuất tài liệu y khoa liên quan.

Pipeline:

\begin{verbatim}
User Query
    |
    v
Embedding
    |
    v
FAISS
    |
    v
Top-K Documents
\end{verbatim}

Ví dụ chunk:

\begin{quote}
Melanoma thường biểu hiện bất đối xứng, nhiều màu sắc và bờ không đều.
\end{quote}

\section{Graph Reasoning}

Mục tiêu:

\begin{verbatim}
Features
   |
   v
Disease Ranking
\end{verbatim}

Ví dụ:

\begin{verbatim}
border_irregularity
color_variation
asymmetry
\end{verbatim}

sẽ làm tăng điểm của MEL.

Công thức:

\[
Score(d)
=
\sum_i
Weight_i \times Confidence_i
\]

\section{Fusion Layer}

Fusion Layer kết hợp nhiều nguồn thông tin.

\subsection{CNN}

\[
P_{cnn}(d)
\]

\subsection{Graph}

\[
P_{graph}(d)
\]

\subsection{RAG Retrieval}

\[
P_{rag}(d)
\]

\subsection{Điểm cuối}

\[
FinalScore(d)
=
w_1P_{cnn}(d)
+
w_2P_{graph}(d)
+
w_3P_{rag}(d)
\]

Trong đó:

\[
w_1+w_2+w_3=1
\]

\section{Explanation Engine}

Mục tiêu là giải thích dự đoán.

Ví dụ:

\begin{verbatim}
Prediction: MEL

Reason:
- Border irregularity detected
- Color variation detected
- CNN confidence high
- Medical evidence supports MEL
\end{verbatim}

\section{Lộ trình mở rộng}

\subsection{Giai đoạn 1}

HAM10000:

\begin{itemize}
\item MEL
\item NV
\item BCC
\item AKIEC
\item BKL
\item DF
\item VASC
\end{itemize}

\subsection{Giai đoạn 2}

Mở rộng lên 20-30 bệnh:

\begin{itemize}
\item Eczema
\item Psoriasis
\item Rosacea
\item Vitiligo
\item Seborrheic Dermatitis
\end{itemize}

\subsection{Giai đoạn 3}

Mở rộng lên 100+ bệnh:

\begin{itemize}
\item Disease Nodes
\item Feature Nodes
\item Symptom Nodes
\item Treatment Nodes
\item Diagnosis Nodes
\end{itemize}

\section{Node Embeddings}

Mỗi node trong graph đều có embedding riêng.

Ví dụ:

\begin{verbatim}
border_irregularity
itching
ulceration
MEL
BCC
\end{verbatim}

đều có vector biểu diễn.

Điều này giúp thực hiện:

\begin{verbatim}
Query
  |
  v
Node Retrieval
  |
  v
Graph Expansion
\end{verbatim}

\section{Pipeline hoàn chỉnh}

\begin{center}
\begin{tikzpicture}[node distance=1.1cm]

\node (a) {Image Input};
\node (b) [below of=a] {CNN};
\node (c) [below of=b] {Feature Layer};
\node (d) [below of=c] {Knowledge Graph};
\node (e) [below of=d] {Vector Retrieval};
\node (f) [below of=e] {Fusion Layer};
\node (g) [below of=f] {Disease Ranking};
\node (h) [below of=g] {Explanation Engine};
\node (i) [below of=h] {Chatbot};

\draw[->] (a)--(b);
\draw[->] (b)--(c);
\draw[->] (c)--(d);
\draw[->] (d)--(e);
\draw[->] (e)--(f);
\draw[->] (f)--(g);
\draw[->] (g)--(h);
\draw[->] (h)--(i);

\end{tikzpicture}
\end{center}

\section{Kiến trúc đích}

Khi hệ thống đạt quy mô hàng trăm bệnh:

\begin{verbatim}
Image
 +
Text Query
        |
        v
Embedding Layer
        |
        v
Knowledge Graph
        +
Node Embeddings
        +
Vector Database
        |
        v
Graph Reasoning
        |
        v
Fusion Layer
        |
        v
Explanation Engine
        |
        v
Chatbot
\end{verbatim}

Đây là kiến trúc phù hợp để mở rộng lâu dài thay vì phụ thuộc hoàn toàn vào rule-based feature extraction.

\end{document}
